\documentclass[11pt,a4paper,reqno,french,tikz]{amsart}
\usepackage[utf8]{inputenc} \usepackage[T1]{fontenc}
%\usepackage{babel} % français \usepackage[latin1]{inputenc}
\usepackage{amsthm,amsmath,amsfonts,amssymb,amsxtra}
\usepackage{appendix,bookmark,comment,cite,hyperref,color,xcolor,cite,graphicx}

%%%%%%%%%%%%%%%%%%%% Standard macros
% Environments
\newtheorem{theorem}{Theorem}[section]\newtheorem{definition}[theorem]{Definition}\newtheorem{lemma}[theorem]{Lemma}\newtheorem{example}[theorem]{Example}\newtheorem{proposition}[theorem]{Proposition}\newtheorem{corollary}[theorem]{Corollary}\newtheorem{conjecture}[theorem]{Conjecture}\newtheorem{remark}[theorem]{Remark}

% Mathbb
\let\C\relax\newcommand{\C}{\mathbb{C}}\newcommand{\Z}{\mathbb{Z}}\newcommand{\R}{\mathbb{R}}\newcommand{\N}{\mathbb{N}}

% Mathcal
\newcommand\cA{\mathcal{A}}\newcommand\cB{\mathcal{B}}\newcommand\cC{\mathcal{C}}\newcommand\cD{\mathcal{D}}\newcommand\cE{\mathcal{E}}\newcommand\cF{\mathcal{F}}\newcommand\cG{\mathcal{G}}\newcommand\cH{\mathcal{H}}\newcommand\cI{\mathcal{I}}\newcommand\cJ{\mathcal{J}}\newcommand\cK{\mathcal{K}}\newcommand\cL{\mathcal{L}}\newcommand\cM{\mathcal{M}}\newcommand\cN{\mathcal{N}}\newcommand\cO{\mathcal{O}}\newcommand\cP{\mathcal{P}}\newcommand\cQ{\mathcal{Q}}\newcommand\cR{\mathcal{R}}\newcommand\cS{\mathcal{S}}\newcommand\cT{\mathcal{T}}\newcommand\cU{\mathcal{U}}\newcommand\cV{\mathcal{V}}\newcommand\cW{\mathcal{W}}\newcommand\cX{\mathcal{X}}\newcommand\cY{\mathcal{Y}}\newcommand\cZ{\mathcal{Z}}

% Math operators
\DeclareMathOperator{\tr}{Tr}\DeclareMathOperator{\ran}{Ran}\DeclareMathOperator{\Span}{Span}\DeclareMathOperator{\Ker}{Ker}\DeclareMathOperator{\re}{Re}\DeclareMathOperator{\id}{id}\DeclareMathOperator{\im}{Im}\DeclareMathOperator{\dist}{dist}\let\div\relax\DeclareMathOperator{\div}{div}
\def\d{{\rm d}}
\usepackage{dsfont} % to define \mathds{1}
\def\1{{\mathds{1}}}

% Others
\renewcommand{\ge}{\geqslant}\renewcommand{\le}{\leqslant}

%%%%%%%%%%%%%%%%%%%% Macros Louis
% Delimiters
\newcommand{\pa}[1]{\left( #1 \right)}                              % ()
\newcommand{\bpa}[1]{\big( #1 \big)}                                % ()
\newcommand{\acs}[1]{\left\{ #1 \right\}}                       % {}
\newcommand{\seg}[1]{\left[ #1 \right]}                             % []
\newcommand{\ab}[1]{\left|#1\right|}                                % ||
\newcommand{\ps}[1]{\left< #1 \right>}                              % <>
\newcommand{\nor}[2]{ \left| \! \left| #1 \right| \! \right|_{#2} } % ||.||_{.}
\newcommand{\norm}[1]{ \left| \! \left| #1 \right| \! \right| }     % ||.||

% Greek letters shortcuts
\newcommand\vp{\varphi}                            % φ
\def\eps{\varepsilon}\newcommand{\ep}{\varepsilon} % ε
\let\p\relax\newcommand{\p}{\psi}                  % ψ
\newcommand{\na}{\nabla}                           % ∇

% Others
\newcommand{\f}[2]{\frac{#1}{#2}}                                              % fraction
\newcommand{\bul}{$\bullet$ \hspace{0.1cm}}                                    % •
\newcommand{\mymax}[1]{\underset{\substack{#1}}{\text{\normalfont{max}}}\;} % max
\newcommand{\mymin}[1]{\underset{\substack{#1}}{\text{\normalfont{min}}}\;} % min
\newcommand{\mysup}[1]{\underset{\substack{#1}}{\text{\normalfont{sup}}}\;} % sup
\newcommand{\myinf}[1]{\underset{\substack{#1}}{\text{\normalfont{inf}}}\;} % inf
\newcommand{\ind}[1]{_{\textup{#1}}}                                           % indice
\newcommand{\mat}[1]{\begin{pmatrix} #1 \end{pmatrix}}                         % matrices
\usepackage{listings}

%%%%%%%%%%%%%%%%%%%% End macros

\title[]{Efficiently compute stuff}
\author[L. Garrigue]{Louis Garrigue}
\address{Laboratoire AGM - UMR 8088, CNRS - CY Cergy Paris Université, 2 avenue Adolphe Chauvin, Cergy-Pontoise, 95302 cedex, France}
\email{louis.garrigue@cyu.fr}
\date{\today}
\begin{document}
\maketitle

There will be
\begin{itemize}
\item object (struct in julia) which contains all the parameters for one computation, type is \texttt{OneComputation} and the objects are denoted by \texttt{oc}
\item object (type \texttt{Stock}) which stores all the solutions, called \texttt{St\_comp}
\end{itemize}
{\color{red}{\texttt{oc} is very light while \texttt{St\_comp} is very heavy}}
Rules :
\begin{itemize}
\item Any quantity can be computed having only one argument \texttt{oc}, possibly using only part of the arguments of \texttt{oc}. It's not a problem that those quantities are not stored because they are quick to compute (for instance the length of an object)
\item Objects of type \texttt{OneComputation} have the minimal amount of variables
\end{itemize}

\subsection{Functions defined in the stock creation}%
\label{sub:Functions defined in the stock creation}

\begin{itemize}
\item The functions defined in the stock creation will take the minimal amount of variables, to minimize the number of times they are computed and stored
\item Define {\color{red}{only}} stored objects which are going to be {\color{red}{long}} to compute. If things are easy to compute (for instance the length of an object), defined them with external functions
\end{itemize}

\begin{lstlisting} 
	function compute_map(coords,St_comp)
		xi, precision = coords
		xi + precision + X("p",xi,St_comp)
	end
\end{lstlisting}


\subsection{External functions}%
\label{sub:External functions}

Those are the functions which are defined outside the stock functions.

They in general take in argument \texttt{oc} and \texttt{St\_comp}. They don't use all the arguments. For instance to get the ``maps'' we can define
\begin{lstlisting} 
function get_map(oc,St_comp)
	coords = oc.xi, oc.precision
	X("map",coords,St_comp)
end
\end{lstlisting}
So each object defined in stock has its external counterpart to call it. But other derivatives can also be defined like
\begin{lstlisting} 
function get_map_length(oc,St_comp)
	map = get_map(oc,St_comp)
	length(map)
end
\end{lstlisting}


\bibliographystyle{siam}
\bibliography{/home/louis/Documents/biblio}
\end{document}
